\subsection{Энергетический подход}
%%%%%%%%%%%%%%%%%%%%%%%%%%%%%%%%%%%%%%%%%%%%%%%%%%%%
%%%%%%%%%%%%%%%%%%%%%%%%%%%%%%%%%%%%%%%%%%%%%%%%%%%%
 \z[{\begin{wrapfigure}[3]{r}{0.3\textwidth}%
		\vspace{-15pt}%
		\begin{flushright}%
 			\includegraphics[width=0.3\textwidth]{10.2.1}%
		\end{flushright}%
	\end{wrapfigure}%
	}]%
	{Пружина жёсткости $k$ одним концом присоединена к оси колеса массы $m$, которое способно катиться без проскальзывания, а другим прикреплена к стенке. Какова частота колебаний системы? Масса колеса однородно распределена по ободу.}
%
{$\omega = \sqrt{\dfrac{k}{2m}}$}
%%%%%%%%%%%%%%%%%%%%%%%%%%%%%%%%%%%%%%%%%%%%%%%%%%%%
\z[{\begin{wrapfigure}[4]{r}{0.3\textwidth}%
		\vspace{-10pt}%
		\begin{flushright}%
 			\includegraphics[width=0.3\textwidth]{10.2.2}%
		\end{flushright}%
	\end{wrapfigure}%
	}]%
	{К ободу колеса с горизонтально расположенной осью прикрепили грузик массы $m$. Найдите массу колеса, предполагая её однородно распределённой по ободу, если частота колебаний колеса с грузиком вокруг оси равна $\omega$, а его радиус равен $R$, $R < \frac{g}{\omega^2}$.}
%
{$M = m \left(\frac{g}{\omega^2 R} - 1\right)$}
%%%%%%%%%%%%%%%%%%%%%%%%%%%%%%%%%%%%%%%%%%%%%%%%%%%% 
 \z{Как изменится частота колебаний маятника, представляющего собой груз на лёгком стержне, если к середине стержня прикрепить горизонтальную пружину жёсткости $k$? На рисунке изображено положение равновесия.}
%
{$\omega^2 = \frac{g}{l} + \frac{k}{4m}$, квадрат частоты возрастет на $\frac{k}{4m}$}
%%%%%%%%%%%%%%%%%%%%%%%%%%%%%%%%%%%%%%%%%%%%%%%%%%%% 
 \z{Два тела массы $m_1$ и $m_2$ связаны пружиной жёсткости $k$. Какова частота свободных колебаний такой системы, если вращения нет?}
%
{$\omega = \sqrt{\frac{k}{\mu}}$, где $\mu = \frac{m_1 m_2}{m_1 + m_2}$ - приведённая масса}
%%%%%%%%%%%%%%%%%%%%%%%%%%%%%%%%%%%%%%%%%%%%%%%%%%%%	
 \z[{\begin{wrapfigure}[3]{r}{0.2\textwidth}%
		\vspace{-15pt}%
		\begin{flushright}%
 			\includegraphics[width=0.2\textwidth]{10.2.5}%
		\end{flushright}%
	\end{wrapfigure}%
	}]%
	{На гладкой горизонтальной поверхности находится тележка массы $M$ с установленным на ней математическим маятником длины $l$ и массы $m$. Найдите период колебаний системы.}
%
{$T = 2\pi \sqrt{\frac{l(m + M)}{Mg}}$}